\par\Needspace{6\baselineskip}\medskip
\SourceLeaf{099}\hypertarget{native-099-article075}{}\pdfbookmark[2]{75.}{article-075}
\begin{center}75.\end{center}
Eine weitere Lösung der vorigen Aufgabe beruht auf der Bemerkung in Artikel 64, III, dass der heliozentrische Ort der Erde, der geozentrische Ort des Himmelskörpers und dessen heliozentrischer Ort auf ein und demselben Großkreis der Kugel liegen. Diese Orte seien in Abb.~3 der Reihe nach $T$, $G$ und $H$. Ferner sei $\AscendingNode$ der Ort des aufsteigenden Knotens, $\AscendingNode T$ und $\AscendingNode H$ seien Teile der Ekliptik und der Bahn, und $GP$ sei das von $G$ auf die Ekliptik gefällte Lot, das daher $=b$ ist. Hieraus und aus dem Bogen $PT=L-l$ bestimmen wir den Winkel $T$ und den Bogen $TG$. Im sphärischen Dreieck $\AscendingNode HT$ sind dann der Winkel $\AscendingNode=i$, der Winkel $T$ und die Seite $\AscendingNode T=L-\AscendingNode$ gegeben. Daraus findet man die beiden übrigen Seiten $\AscendingNode H=u$ und $TH$. Schließlich gilt $HG=TG-TH$ sowie $r=\frac{R\sin TG}{\sin HG}$, $\Delta=\frac{R\sin TH}{\sin HG}\mathpunct{\cdot}$
